% !PRINTED 166
\weberpage{166}

%<<CONT
zwei Functionen $UV$ nur dann durch $v$ theilbar sein kann, wenn
einer der beiden Factoren durch $v$ theilbar ist. Dasselbe gilt
für ein Product von mehreren Functionen, und so folgt aus I.
das Theorem:

\gesperrt{II. Ein Product aus mehreren ganzen rationalen Functionen ist nur dann durch eine unzerlegbare Function $v$ theilbar, wenn wenigstens einer der Factoren des Productes durch $v$ theilbar ist.}

Wenn eine ganze rationale Function $U$ auf zwei Arten in
unzerlegbare Factoren zerlegt ist,
\[
U = v\,v'\,v''\dots = w\,w'\,w''\dots,
\]
so muss nach II. wenigstens einer der Factoren $v, v', v''\dots$
durch $w$ theilbar sein, also etwa $v$. Dann aber kann, da auch
$v$ unzerlegbar ist, $v$ von $w$ nur durch einen constanten Factor
verschieden sein.

Demnach ist, wenn $c$ dieser constante Factor ist,
\[
c\,v'\,v''\cdots = w'\,w''\dots,
\]
woraus folgt, dass eine der Functionen $v', v''\dots$, etwa $v'$, durch
$w'$ theilbar ist, und sich also von $w'$ nur durch einen constanten
Factor unterscheidet.

Wir erhalten also als zweite Folgerung aus dem Theorem I:

\gesperrt{III. Eine ganze rationale Function kann, von constanten Factoren abgesehen, nur auf eine Art in unzerlegbare Factoren zerlegt werden.}

% [corrected 2026-08-07, 8x zoom strip 6/8] "grössten gemeinschaftlichen Theilers" is
% letterspaced -- the term at its definition. "des" before it and "von zwei oder mehr" after it
% are set normally, so the run is exactly those three words. Raised by the sperrung pass; the
% numbered enunciation III above it was already correctly wrapped, which is what makes the
% contrast on this page easy to read at zoom.
Hieraus ergiebt sich der Begriff des \gesperrt{grössten gemeinschaftlichen
Theilers} von zwei oder mehr ganzen rationalen Functionen
$U, V\dots$ Man versteht darunter das Product aller
unzerlegbaren Factoren, die in den Zerlegungen jeder der Functionen
$U, V\dots$ vorkommen, oder die Function möglichst hohen
Grades, die in allen Functionen $U, V\dots$ aufgeht. Nach III. ist
diese Function, von einem constanten Factor abgesehen, für jedes
Functionensystem $U, V\dots$ vollständig bestimmt.

Mehrere Functionen $U, V, W\dots$ heissen relativ prim, wenn
keine zwei von ihnen einen gemeinsamen Theiler haben.

Alle diese Definitionen und Sätze sind genau analog mit
sehr bekannten Sätzen der elementaren Zahlenlehre, die dort als
Folgerungen des Algorithmus des grössten gemeinschaftlichen
%CONT>>
